\documentclass[11pt]{article}

\usepackage[margin=1in]{geometry}
\usepackage[utf8]{inputenc}
\usepackage[T1]{fontenc}
\usepackage{amsmath,amssymb,amsthm,bm}
\usepackage{booktabs}
\usepackage{array}
\usepackage{graphicx}
\usepackage{tikz}
\usepackage{pgfplots}
\pgfplotsset{compat=1.17}
\usetikzlibrary{positioning,arrows.meta,calc}
\usepackage{caption}
\captionsetup{font=small,labelfont=bf,margin=0.5cm}
\usepackage[round,authoryear]{natbib}
\usepackage{hyperref}
\usepackage{microtype}
\hypersetup{colorlinks=true,linkcolor=black,citecolor=black,urlcolor=blue}

\usepackage{xcolor}
\definecolor{iconblue}{RGB}{30,72,116}
\definecolor{iconaccent}{RGB}{198,108,42}
\definecolor{icongray}{RGB}{90,90,90}

\newcommand{\R}{\mathbb{R}}
\newcommand{\E}{\mathbb{E}}
\newcommand{\Tth}{T_{\!\theta}}
\newcommand{\norm}[1]{\left\lVert #1 \right\rVert}
\DeclareMathOperator{\bool}{bool}

\IfFileExists{toy_multiphysics/metrics.tex}{\input{toy_multiphysics/metrics.tex}}{}
\IfFileExists{toy_multiphysics/eval_metrics.tex}{\input{toy_multiphysics/eval_metrics.tex}}{}
\providecommand{\ToyGrid}{64}
\providecommand{\ToyFrames}{4}
\providecommand{\ToyShots}{4}
\providecommand{\ToyTrainTime}{$\sim\!30$~min}
\providecommand{\ToyTestRel}{15.8\%}
\providecommand{\ToyContextSummary}{0-shot 22.0\%, 3-shot 15.0\%}
\providecommand{\ToyInverseAcc}{81.0\%}
\providecommand{\ToyZeroShotImp}{16.4\%}

\title{\vspace{-1.5em}
  ICON-LM for In-Context Viscous Burgers Solvers:\\
  A Reproducible Implementation Study}

\author{
  Aditya Patil, Dhruv Patel\\
  \small MATH 270A Final Report\\
  \small June 9, 2026
}
\date{}

\begin{document}
\maketitle
\vspace{-2em}

%==============================================================
\begin{abstract}
\noindent
Operator learning asks a neural network to approximate maps between functions, such
as the map from an initial condition to the solution of a partial differential equation (PDE). ICON-LM extends this idea by making the examples themselves a prompt. In it,
a transformer reads several condition--solution pairs, an optional natural-language
caption, and uses that to predict a new solution in a single forward pass. This report 
summarizes the mathematical structure of ICON-LM and then documents our
implementation of a smaller ICON-LM-style model used on the one-dimensional viscous
Burgers equation. In our experiment, each operator is indexed by an unknown
viscosity $\nu$, and the model must infer the appropriate dynamics from sparse
trajectory examples and optional captions. The code generates spectral-solver data,
builds the ICON-LM attention mask, trains a 2.21M-parameter transformer, and
evaluates context scaling from zero to three in-context demonstrations. On the
saved test set, no-caption relative error improves from $23\%$ at zero shots to
$16\%$ at three shots, while captions help reduce the zero-shot error. These results
support the intuition behind ICON-LM, but also demonstrate limitations such as how caption embeddings need very careful numerical handling, and that small models only weakly use exact parameter information.
\end{abstract}

%==============================================================
\section{Introduction}
\label{sec:intro}

Many scientific computing problems can be written as operator-learning problems.
Instead of predicting a scalar label from a vector input, the model predicts an
entire function from another function. For a time-dependent PDE, the input might be
an initial condition $u_0(x)$ and the output might be the future trajectory
$u(x,t)$. Classical neural operator methods such as DeepONet and the Fourier
Neural Operator learn this map after training on many examples from a fixed
operator~\citep{lu2021deeponet,li2021fourier,kovachki2023neural}. They lower
the cost of solving one PDE family, but a change in coefficients, boundary
conditions, or physical parameters requires retraining or fine-tuning.

In-Context Operator Networks (ICON) propose a different abstraction, where the neural
network is trained not merely to approximate one operator, but to infer an
operator from examples supplied at inference time~\citep{yang2023icon}. Given a
prompt of condition--quantity-of-interest pairs
$\{(c_j,u_j)\}_{j=1}^J$ and a new condition $c_{\rm q}$, the model predicts
$u_{\rm q}$ without any gradient updates. This is directly analogous to in-context
learning in language models~\citep{brown2020language}, just with functions replacing
sentences and operators replacing tasks.

ICON-LM strengthens this analogy~\citep{yang2024icon_lm}. It replaces the
encoder--decoder ICON architecture with a single transformer encoder and a
carefully designed attention mask, which makes the model behave like an
autoregressive learner over functions. Earlier examples can inform later
predictions, but the answer tokens for a current example can't leak into that
example's query tokens. ICON-LM also adds an optional caption modality, allowing the prompt
to contain natural language text such as the governing equation or numerical parameter values.

Our project has two parts. First, we analyzed the ICON-LM formulation, as it's
the conceptual basis for our work. Second, we implemented a compact, reproducible
version for a single PDE family, the 1-dimensional
viscous Burgers Equation. The implementation is deliberately smaller than the
paper model, but it still utilizes the governing PDE, a numerical data generator, function tokenization, the
ICON-LM mask, next-function training, caption conditioning, and quantitative
evaluation.

\paragraph{Code availability:} 
The code implementation used for this report is available at \\
\href{https://github.com/dhruvpatel-05/270A-finalproject}{https://github.com/dhruvpatel-05/270A-finalproject}.

%==============================================================
\section{Mathematical and Scientific Formulation}
\label{sec:formulation}

\subsection{In-Context Operator Learning}

Let $\mathcal{C}$ and $\mathcal{U}$ be spaces of input and output functions. An
operator is a map
\begin{equation}
T:\mathcal{C}\to\mathcal{U}, \qquad c\mapsto u=T(c).
\end{equation}
In a standard operator-learning setup, a model is trained for one fixed $T$. In an
in-context setup, there is a distribution $\mathcal{D}$ over operators. At
inference, the model receives a prompt
\begin{equation}
P_J=\left(\{(c_j,u_j)\}_{j=1}^{J},c_{\rm q}\right),
\end{equation}
where all pairs are generated by the same but unknown operator $T\sim\mathcal{D}$.
The learned model $\Tth$ produces $\hat u_{\rm q}=\Tth(P_J)$. A natural meta-loss is
\begin{equation}
\mathcal{L}(\theta)=
\E_{T\sim\mathcal{D}}\,
\E_{P_J}\,
\E_{y}
\left[
  \left|\Tth(P_J)(y)-T(c_{\rm q})(y)\right|^2
\right],
\label{eq:meta_loss}
\end{equation}
where $y$ is an evaluation point for the output function. The important distinction
is that $T$ is not encoded in $\theta$ alone, as it must be inferred from the prompt.

The multi-modal extension adds a caption $\kappa$, so the model becomes
$\Tth(\kappa,P_J)$. In the ICON-LM paper, vague captions describe the equation
class while precise captions include parameter values~\citep{yang2024icon_lm}. In
our experiment, precise captions include the Burgers viscosity, while vague
captions state only whether the flow is low-viscosity and sharp or high-viscosity
and smooth.

\subsection{Burgers Operator Family}

The implemented scientific system is the one-dimensional viscous Burgers equation
on the periodic domain $[0,2\pi)$:
\begin{equation}
u_t + u u_x = \nu u_{xx}, \qquad x\in[0,2\pi),\quad t\in[0,1].
\label{eq:burgers}
\end{equation}
The hidden operator parameter is the viscosity
\begin{equation}
\nu \in [10^{-2},10^{-1}].
\end{equation}
For each fixed $\nu$, the operator maps an initial condition to a future
trajectory,
\begin{equation}
T_{\nu}: u_0(x)\mapsto \left(u(x,t_1),\ldots,u(x,t_K)\right),
\qquad t_k\in\{0.25,0.50,0.75,1.00\}.
\end{equation}
Low viscosity produces sharper gradients and more advection-dominated behavior;
high viscosity damps high-frequency content and gives smoother solutions. This
makes the family scientifically meaningful while still small enough to train
locally.

For our initial conditions, we use random Fourier sums:
\begin{equation}
u_0(x)=\frac{1}{5}\sum_{m=1}^{5} a_m\sin(mx+\phi_m),
\qquad
a_m\sim{\rm Unif}[-1,1],\quad
\phi_m\sim{\rm Unif}[0,2\pi].
\label{eq:ic}
\end{equation}
The data generator uses $N=64$ uniform grid points and solves
Equation~\eqref{eq:burgers} with a Fourier spectral right-hand side and SciPy
\texttt{RK45} time stepping. Spatial derivatives are computed in Fourier space:
\begin{equation}
\widehat{u_x}_k = i k \hat u_k,
\qquad
\widehat{u_{xx}}_k = -k^2 \hat u_k.
\end{equation}
The saved dataset contains $200$ training viscosities on a log-spaced grid and
$40$ held-out viscosities sampled log-uniformly, with $8$ initial conditions per
viscosity.

\subsection{Evaluation Metrics}

For trajectory prediction, we report relative $L_2$ error over all query tokens:
\begin{equation}
{\rm RelErr} =
100\left(
\frac{\sum_{b,q} \left(\hat u_{bq}-u_{bq}\right)^2}
     {\sum_{b,q} u_{bq}^2}
\right)^{1/2}.
\label{eq:relerr}
\end{equation}
This is computed after undoing training normalization. We also
include a diagnostic inverse task, where a nearest-centroid classifier estimates which
log-viscosity bin generated a trajectory. That inverse classifier isn't the
ICON-LM model, but serves as a check that the dataset contains physically recoverable
information about $\nu$.

%==============================================================
\section{Methodology and Implementation}
\label{sec:methodology}

\subsection{Function Tokens}

The original ICON-LM represents scattered function samples as tokens. Our
implementation follows that principle with three token types for each trajectory:
condition tokens $C$, observed future-solution tokens $U$, and query tokens $Q$.
Each token has the form
\begin{equation}
(\tau,x,t;v),
\end{equation}
where $\tau\in\{C,U,Q\}$ is the token type, $(x,t)$ is the coordinate, and $v$ is
the scalar value supplied to the transformer.

For each source trajectory, the condition block contains $n_C=32$ sparse samples
of the initial condition at $t=0$:
\begin{equation}
C_j = \{(C,x_i,0;u_{0,j}(x_i))\}_{i=1}^{32}.
\end{equation}
The solution block contains $32$ spatial samples at each of the four future times,
yielding $128$ observed future tokens:
\begin{equation}
U_j = \{(U,x_i,t_k;u_j(x_i,t_k))\}_{i,k}.
\end{equation}
The query block has the same future coordinates, but its input value is the local
initial condition value rather than the target future value:
\begin{equation}
Q_j = \{(Q,x_i,t_k;u_{0,j}(x_i))\}_{i,k},
\qquad
\text{target } u_j(x_i,t_k).
\end{equation}
Thus the model is told where to predict and what the initial condition is at that
location, but it is never given the future answer for a current query token.

\subsection{Attention Mask}

The core technical device is the ICON-LM attention mask. Let token $a$ be the row
token asking for information and token $b$ be the column token being attended to.
Let $e(a)$ be the example index and $\tau(a)$ be the type. In our no-caption mask,
attention is allowed when any one of the following holds:
\begin{align}
e(b)<e(a) &\quad\text{and}\quad \tau(b)\in\{C,U\}, \\
e(b)=e(a),\ \tau(a)=C &\quad\text{and}\quad \tau(b)=C, \\
e(b)=e(a),\ \tau(a)=U &\quad\text{and}\quad \tau(b)\in\{C,U\}, \\
e(b)=e(a),\ \tau(a)=Q &\quad\text{and}\quad \tau(b)=C, \\
a=b,\ \tau(a)=Q. &&
\end{align}
This gives the exact behavior needed for next-function prediction. Query tokens
can use earlier condition--solution examples and the current condition, but can't
see the current solution tokens $U_j$. Query tokens also cannot attend to other
query tokens, so predictions at different coordinates do not communicate through
hidden answer-dependent channels. As in the ICON-LM paper, this prevents both
direct leakage and multi-layer indirect leakage.

Figure~\ref{fig:mask} shows the full captioned mask used to explain this rule. The actual code constructs the corresponding Boolean matrix and includes tests to check that current-query tokens
can't attend to current-solution tokens and that query-query attention is
diagonal.

\begin{figure}[t]
\centering
\IfFileExists{mask.png}{%
  \includegraphics[width=0.98\textwidth]{mask.png}%
}{%
  \fbox{\parbox{0.92\textwidth}{\centering Upload \texttt{mask.png} to render
  the captioned attention-mask figure.}}%
}
\caption{Captioned ICON-LM attention mask. Rows are querying tokens and columns
are attended tokens. The caption block is causal (Previous condition/QoI tokens
are visible to later examples), and each query token can see the caption, previous
examples, the current condition, and itself, but not the current QoI answer.}
\label{fig:mask}
\end{figure}

\subsection{Model and Captions}

The model is a pre-layer-norm transformer encoder written in PyTorch. Each token
embedding is the sum of a type embedding, a linear projection of $(x,t)$, a linear
projection of the scalar value, and a learned demonstration-position embedding
$3j+\tau$. The saved configuration is summarized in Table~\ref{tab:impl_config}.

\begin{table}[t]
\centering
\small
\begin{tabular}{ll}
\toprule
Component & Value in saved run \\
\midrule
PDE grid and times & $N=64$, $K=4$, $t=(0.25,0.50,0.75,1.00)$ \\
Training/test operators & $200$ train viscosities, $40$ held-out viscosities \\
Sources per operator & $8$ initial conditions \\
Prompt layout & $4$ training examples, $32$ condition samples, $128$ future samples \\
Transformer & $d=256$, $4$ layers, $8$ heads, feed-forward width $512$ \\
Parameters & $2.21$M \\
Optimizer & AdamW, learning rate $3\times10^{-4}$, cosine decay, gradient clip $1.0$ \\
Caption use during training & $60\%$ caption batches; half precise, half vague \\
Saved evaluation & Relative $L_2$ on held-out viscosities and initial conditions \\
\bottomrule
\end{tabular}
\caption{Implementation details.}
\label{tab:impl_config}
\end{table}

Captions are represented differently from the full ICON-LM paper. Instead of
fine-tuning GPT-2, the code uses a frozen MiniLM sentence
encoder~\citep{reimers2019sentence} to embed one caption string into a vector. That
vector is projected to the transformer width and prepended as a single prefix
token. Caption attention is prefix-style, so the physics tokens attend to the caption
token, while the caption token doesn't attend back to the physics sequence.

The precise caption includes text such as
\emph{``Burgers PDE with viscosity nu=0.02500''}. However, in practice, sentence embeddings
are weak at distinguishing nearby numeric strings. Vague captions only indicate the physical regime, for example ``strong viscous
damping'' or ``weak viscosity and sharp gradients.''

\subsection{Training Objective}

Training uses the ICON-LM next-function idea. A batch samples one viscosity and
four source trajectories for each prompt. For every example $j$, the query block
tries to reconstruct the corresponding future solution. With a caption, the first
example is included in the loss, so zero-shot behavior is trained. Without a
caption, the first example is excluded as the model has neither prior examples
nor an operator description. In normalized units, the loss equation is
\begin{equation}
\mathcal{L}_{\rm prompt}(\theta)=
\frac{1}{|\mathcal{Q}|}
\sum_{j\in\mathcal{J}_{\rm loss}}
\sum_{(x,t)\in Q_j}
\left(
\hat u_{\theta,j}(x,t)-u_j(x,t)
\right)^2,
\label{eq:toy_loss}
\end{equation}
where $\mathcal{J}_{\rm loss}=\{1,2,3,4\}$ for prompts with captions and
$\mathcal{J}_{\rm loss}=\{2,3,4\}$ for prompts without captions.

At test time, the evaluation script builds prompts with $J_{\rm demo}$ completed
demonstration trajectories and one final query trajectory. We evaluate
$J_{\rm demo}\in\{0,1,2,3\}$ because this keeps the final query at demonstration
positions that were directly trained as targets.

%==============================================================
\section{Results and Analysis}
\label{sec:results}

\subsection{Forward Prediction and Context Scaling}

Table~\ref{tab:context} reports the saved sparse-query evaluation on the held-out
test set. The main trend matches the ICON-LM hypothesis that examples in the prompt
help the model specialize to the hidden operator. With no caption, relative error
falls from $22.77\%$ at zero demonstrations to $15.95\%$ at three demonstrations.
Captions mostly help when the prompt has few or no examples. At zero shots, the
vague-caption model improves over the no-caption baseline by $6.62$ percentage points,
and the precise-caption model improves by $8.22$ percentage points.

\begin{table}[t]
\centering
\small
\begin{tabular}{lcccc}
\toprule
Prompt mode & 0-shot & 1-shot & 2-shot & 3-shot \\
\midrule
Precise caption & $14.55\pm0.37\%$ & $13.92\pm0.41\%$ & $14.29\pm0.26\%$ & $14.38\pm0.42\%$ \\
Vague caption   & $16.15\pm0.62\%$ & $14.70\pm0.47\%$ & $15.09\pm0.29\%$ & $15.06\pm0.49\%$ \\
No caption      & $22.77\pm0.35\%$ & $15.61\pm0.57\%$ & $16.08\pm0.38\%$ & $15.95\pm0.59\%$ \\
\bottomrule
\end{tabular}
\caption{Relative $L_2$ prediction error on held-out
Burgers operators. Errors are mean $\pm$ std over three seeds.}
\label{tab:context}
\end{table}
The same context-scaling result is visualized in Figure~\ref{fig:context_scaling}.
The plot makes the main effect clearer than the table: the no-caption model is
poor at zero shots because it has no information about the hidden viscosity, while
one in-context trajectory closes most of the gap.
\begin{figure}[t]
\centering
\includegraphics[width=0.64\textwidth]{panel_a.png}
\caption{Few-shot context scaling on the one-dimensional Burgers experiment.
This compact plot corresponds to 
Table~\ref{tab:context}.}
\label{fig:context_scaling}
\end{figure}

The strongest improvement occurs from zero to one demonstration. This is physically
reasonable, as a single observed trajectory reveals whether the viscosity
is damping high-frequency structure quickly or allows sharp gradients to
persist. Additional examples help less because all initial conditions are drawn
from the same random Fourier distribution, so the prompt
already has a fairly clear signal about the operator after just one or two examples.

\subsection{Inverse Diagnostic}

As a separate diagnostic, the evaluation code bins held-out operators by
$\log\nu$ and classifies each operator from five trajectory features: relative
change from the initial condition, final-frame change, frame-to-frame change,
high-frequency power ratio, and gradient magnitude. A nearest-centroid classifier
achieves $\ToyInverseAcc$ accuracy over five bins. Most mistakes occur between
neighboring viscosity bins, which is consistent with the fact that Burgers
dynamics vary continuously with $\nu$. This diagnostic supports the premise that
the context trajectories carry recoverable information about the hidden operator.

%==============================================================
\section{Discussion}
\label{sec:discussion}

\subsection{What Worked}

The experiment successfully implements the central ICON-LM mechanism in a compact
setting. The mask is the most important piece, letting a standard self-attention
stack act as a causal learner over functions while preserving permutation
symmetry inside each sampled function. The same model can run with or without a
caption because captions are just prefix tokens. The code also supports
grid-flexible inference: changing the query-token set is enough to ask for a
different output grid, since the final layer is applied tokenwise.

The results also show meaningful in-context adaptation. The no-caption
zero-shot case is effectively an average-operator baseline, as it sees the initial
condition but has no demonstration revealing the viscosity. Adding one
condition--trajectory example sharply lowers error, indicating that the model uses
the prompt to infer the operator parameter.

\subsection{Limitations and Sources of Error}

The main limitation is scale. The original ICON-LM paper trains on $19$ operator
families and uses much larger models, including a GPT-2 backbone for the
multi-modal setting~\citep{yang2024icon_lm}. Our implementation trains a 2.21M
parameter transformer on one PDE family. It is therefore a faithful mechanism
study, not evidence that the model is a general scientific foundation model.

The caption channel is also simplified. Using MiniLM embeddings avoids the compute
cost of fine-tuning a language model, but it weakens the claim that the model is
parsing equations as language.

There are also numerical and modeling simplifications. The data live on a fixed
periodic one-dimensional grid, the initial conditions come from a narrow Fourier
distribution, and all trajectories are short. The model predicts sampled future
states directly rather than enforcing the PDE residual. Consequently, it can make
physically inconsistent predictions, especially in low-viscosity regimes where
small phase errors near sharp fronts matter. A physics-informed loss or a
conservation-aware architecture could reduce this error, but would move away from
the clean prompt-learning abstraction.

\subsection{Future Improvements}

The most natural next step is to broaden the operator distribution. Even within
Burgers dynamics, one could vary forcing, boundary conditions, final time, or the
initial-condition spectrum. A second direction is to replace the frozen sentence
embedding with a trainable text encoder or a small language-model backbone, then
test whether exact numerical captions become more useful. Finally, evaluation
should include out-of-distribution viscosities and higher-resolution query grids
to measure whether the model has learned an actual robust operator representation rather
than interpolation over the training setup.

%==============================================================
\section{Conclusion}
\label{sec:conclusion}

ICON-LM reframes operator learning as a language-model-like prediction problem
over functions. Our implementation demonstrates this idea on the viscous Burgers
equation, where sparse function samples become tokens, examples become context, captions
become prefix information, and a custom attention mask prevents leakage from
answers into queries. The saved experiment shows that in-context examples reduce
held-out trajectory error, and captions help most when no demonstrations are
available. At the same time, the project makes clear that multi-modal scientific
operator learning is sensitive to representation choices. Exact numerical
information in text isn't automatically useful. The model, caption encoder, and
training distribution must make that information accessible.

%==============================================================
\begin{thebibliography}{99}

\bibitem[Bommasani et~al.(2021)]{bommasani2021opportunities}
R.~Bommasani et~al.
\newblock On the opportunities and risks of foundation models.
\newblock {\em arXiv preprint arXiv:2108.07258}, 2021.

\bibitem[Brown et~al.(2020)]{brown2020language}
T.~Brown et~al.
\newblock Language models are few-shot learners.
\newblock {\em Advances in Neural Information Processing Systems}, 33:1877--1901, 2020.

\bibitem[Karniadakis et~al.(2021)]{karniadakis2021pinn}
G.~E. Karniadakis, I.~G. Kevrekidis, L.~Lu, P.~Perdikaris, S.~Wang, and L.~Yang.
\newblock Physics-informed machine learning.
\newblock {\em Nature Reviews Physics}, 3:422--440, 2021.

\bibitem[Kovachki et~al.(2023)]{kovachki2023neural}
N.~Kovachki, Z.~Li, B.~Liu, K.~Azizzadenesheli, K.~Bhattacharya, A.~Stuart, and A.~Anandkumar.
\newblock Neural operator: Learning maps between function spaces with applications to PDEs.
\newblock {\em Journal of Machine Learning Research}, 24:1--97, 2023.

\bibitem[Li et~al.(2021)]{li2021fourier}
Z.~Li, N.~B. Kovachki, K.~Azizzadenesheli, B.~Liu, K.~Bhattacharya, A.~Stuart, and A.~Anandkumar.
\newblock Fourier neural operator for parametric partial differential equations.
\newblock In {\em International Conference on Learning Representations}, 2021.

\bibitem[Lu et~al.(2021)]{lu2021deeponet}
L.~Lu, P.~Jin, G.~Pang, Z.~Zhang, and G.~E. Karniadakis.
\newblock Learning nonlinear operators via DeepONet based on the universal approximation theorem of operators.
\newblock {\em Nature Machine Intelligence}, 3(3):218--229, 2021.

\bibitem[Reimers and Gurevych(2019)]{reimers2019sentence}
N.~Reimers and I.~Gurevych.
\newblock Sentence-BERT: Sentence embeddings using Siamese BERT-networks.
\newblock In {\em Proceedings of EMNLP-IJCNLP}, 2019.

\bibitem[Vaswani et~al.(2017)]{vaswani2017attention}
A.~Vaswani et~al.
\newblock Attention is all you need.
\newblock In {\em Advances in Neural Information Processing Systems}, 2017.

\bibitem[Yang et~al.(2023)]{yang2023icon}
L.~Yang, S.~Liu, T.~Meng, and S.~J. Osher.
\newblock In-context operator learning with data prompts for differential equation problems.
\newblock {\em Proceedings of the National Academy of Sciences}, 120(39):e2310142120, 2023.

\bibitem[Yang et~al.(2024)]{yang2024icon_lm}
L.~Yang, S.~Liu, and S.~J. Osher.
\newblock Fine-tune language models as multi-modal differential equation solvers.
\newblock Preprint, UCLA, 2024.

\end{thebibliography}

\end{document}
